\section{Evaluation Methodology}
\label{sec:evaluation}

\subsection{Primary and secondary metrics}
Each completed run is scored with the project's deterministic ISV evaluation stack implemented in \texttt{src/isv\_eval/} (CLI entry via the project's \texttt{isv-eval} tooling).
Scoring does not call an LLM.
The metrics used in Phase~2B are:

\begin{itemize}[leftmargin=*]
  \item \textbf{canonical coverage} (primary): share of lexical tokens classified as resource-supported under the canonical inventory policy;
  \item \textbf{broader resource-supported coverage}: canonical-supported tokens plus additional lexical tokens with qualifying exact attestation in audited alternative resources;
  \item \textbf{unresolved rate}: share of lexical tokens in the unresolved (unsupported) bucket;
  \item \textbf{orthography outside-inventory}: character-level count of letters/symbols outside the accepted Interslavic Latin inventory (separate diagnostic in \texttt{src/isv\_eval/orthography.py}).
\end{itemize}

Manuscript quantitative claims focus on canonical coverage and its paired $\Delta$.
Other metrics are retained in the machine-readable datasets for audit and are not collapsed into a single composite score.
In particular, unresolved rate tends to move inversely with canonical coverage under the evaluator construction; reporting both without inventing a composite avoids double-counting the same inventory phenomenon under two names.

\subsection{Operational definition of canonical coverage}
The evaluator tokenizes model output and distinguishes \emph{lexical} (word-like) tokens from non-lexical tokens (punctuation, numbers, and similar).
Coverage denominators use lexical tokens only, as documented in the metrics module denominator policy.

\paragraph{Token and lookup policy.}
A token is treated as lexical when it contains at least one Unicode letter (\texttt{src/isv\_eval/normalize.py}).
Surface forms are not rewritten for display; normalization builds lookup keys only (NFC + lowercase; an optional secondary folded key for selected etymological characters).
Cyrillic tokens may be transliterated to Latin for lookup by the morphology backend.
Each lexical token is classified into inventory buckets used by the metrics module (\texttt{src/isv\_eval/metrics.py}, \texttt{classifier.py}, \texttt{lexicon.py}):
bucket~A (exact dictionary match), bucket~B (morphologically justified under available ISV resources), or bucket~C (unresolved).
Canonical coverage is $(A+B)/\text{lexical\_tokens}$.
Broader coverage additionally counts selected unresolved tokens that carry exact alternative-resource attestation; the broader tier is an evidence estimate, not a validity claim.

\paragraph{Orthography diagnostic.}
Orthography outside-inventory counts characters outside the accepted Interslavic Latin inventory (standard letters, etymological extensions, and sanctioned alternatives documented in the orthography module).
Polish-specific letters absent from that inventory (e.g.\ \emph{\k{a}}, \emph{\l{}}, \emph{\'{o}}, \emph{\.{z}}) are flagged when present.
The orthography check never repairs text and never enters the coverage denominators.

\paragraph{Implementation location.}
Metric definitions and denominator policy are implemented in \texttt{src/isv\_eval/metrics.py}; character audit in \texttt{src/isv\_eval/orthography.py}; Phase~2B aggregates regenerate from per-run \texttt{evaluation.json} / \texttt{orthography.json} via \texttt{scripts/analyze\_exp004\_phase2b.py}.

\subsection{What coverage is not}
Canonical coverage is a screening metric for resource-supported ISV lexical realization.
It is \emph{not} claimed to be equivalent to:
\begin{itemize}[leftmargin=*]
  \item human-judged translation quality,
  \item adequacy/fluency in the sense of DA/MQM protocols,
  \item BLEU~\citep{papineni-etal-2002-bleu} or COMET~\citep{rei-etal-2020-comet}.
\end{itemize}
Orthographic cleanup or substitution of inventory-supported synonyms can move coverage without implying a global quality improvement.
Conversely, communicative improvements that do not increase inventory hits may be invisible to the metric.
This asymmetry is why the paper pairs every major quantitative claim with qualitative audits and explicit limitations.

\subsection{Paired analysis unit}
The experimental unit for priming effects is the Direct$\leftrightarrow$Primed pair within a configuration and replicate, not an unpaired mixture of runs.
Configuration summaries average three paired $\Delta_i$ values.
Invalid pairs are excluded from primary paired inference rather than imputed.

\subsection{Qualitative audit}
For each regime, all valid Direct$\leftrightarrow$Primed pairs were inspected in a descriptive qualitative audit (HIGH and UNSEEN: 21 pairs; LOW primary: 18 valid pairs).
Audits record apparent classes such as lexical borrowing, orthographic cleanup, corpus-opening alignment, and degradations, without adjudicating full linguistic correctness.
Opening-overlap checks against a corpus narrative excerpt are used as an operational detector for winter-opening near-copy, especially to contrast HIGH against LOW/UNSEEN.

\subsection{Reporting conventions used in Results}
Unless otherwise stated, coverage means are percentages averaged over three Direct or three Primed runs, and $\Delta$ means are averages of three paired differences.
We report two decimal places in percentage points when that precision is present in the combined tables.
Sign consistency is reported as all-positive, all-negative, or mixed within the three replicates.
Invalid cells are marked explicitly rather than imputed.
